\section{Conclusion}

This study investigated how various training configurations and model complexities affect the performance of object detection models designed for cow detection in indoor farm environments. Our findings highlight that model performance is highly dependent on camera viewpoints, with side views posing the most significant challenges. We also found that fine-tuning models with weights from similar datasets substantially improves performance, especially for complex models operating with limited data. To further support the research community, we've introduced COLO, a new public dataset for cow localization.

Our results indicate that while increasing model complexity can enhance performance, this isn't universally true, particularly in challenging scenarios, such as training models on top-view images and using them to predict side-view images. Models trained on a single viewpoint demonstrated limited generalization, emphasizing the critical need for diverse and consistent camera angles within training data.

Furthermore, while transfer learning from similar datasets can boost model performance when target datasets are small, it's crucial to ensure the training data accurately represents the target environment. The benefits of transfer learning diminish considerably when the new deployment environment considerably differs from the original training data. We also observed that larger models tend to leverage transfer learning more effectively than smaller models, likely due to their greater capacity to learn complex features.

Despite these valuable insights, our study has certain limitations. The models' performance was evaluated under specific indoor farm conditions, meaning the findings may not directly generalize to all livestock environments. Additionally, our reliance on predefined configurations might limit the applicability of our results to more dynamic settings. Future research should focus on developing adaptive methods to enhance model generalization across various viewpoints and environmental conditions. Specifically, incorporating advanced data augmentation techniques that simulate diverse camera angles in 3D space could effectively address performance degradation caused by different viewpoints.

In conclusion, this study offers practical guidance for reproducing model performance in new environmental settings and contributes the public COLO dataset to foster further research and advancements in this field..
